\chapter{วิเคราะห์ผลการดำเนินงาน}

\section{จุดแข็ง}
\begin{enumerate}[leftmargin=1.6em,itemsep=3pt]
  \item \textbf{ภารกิจหลักตามพระราชดำริทำได้เกินเป้าแล้ว} --- ตัวชี้วัดที่ 39 อยู่ที่ร้อยละ~103
        (62 จากเป้า 60) กระจายใน 37~โครงการ สะท้อนว่างานที่ทำตอบโจทย์พื้นที่จริง
        มิใช่การรับตัวชี้วัดมากระจุกที่โครงการใดโครงการหนึ่ง
  \item \textbf{การกระจายภารกิจครอบคลุมทั้งมหาวิทยาลัย} --- 11~หน่วยงาน ทั้งคณะ สถาบัน วิทยาลัย
        และวิทยาเขตเชียงราย น่าน ตาก พิษณุโลก ทำให้แผนงานไม่ผูกกับหน่วยงานเดียว
  \item \textbf{กลไกอนุมัติและโอนงบประมาณทำงานได้ดี} --- อนุมัติแล้ว 59 จาก 61~โครงการ
        โอนงบถึงผู้รับผิดชอบร้อยละ~98 แสดงว่าคอขวดไม่ได้อยู่ที่ต้นทางของกระบวนการ
  \item \textbf{มีระบบติดตามกลางแล้ว} --- ฐานข้อมูลเชื่อมรหัส ERP เข้ากับโครงการ แผนงาน หน่วยงาน
        และตัวชี้วัด พร้อมแดชบอร์ดสำหรับผู้บริหาร ทำให้รายงานฉบับนี้ตรวจสอบย้อนกลับได้ทุกตัวเลข
\end{enumerate}

\section{ช่องว่างและความเสี่ยง}

\subsection{ความเสี่ยงด้านการเบิกจ่าย (ระดับสูง)}
เหลือเวลาประมาณ 55~วัน แต่ต้องเบิกจ่ายอีก 3,526,251~บาท โดยมีภาระผูกพันใบสั่งซื้อรองรับแล้ว
เพียง 825,156~บาท หมายความว่า\textbf{วงเงินกว่า 2.7 ล้านบาทยังไม่เริ่มกระบวนการจัดซื้อจัดจ้าง}
ซึ่งเป็นข้อจำกัดสำคัญเพราะกระบวนการพัสดุมีระยะเวลาดำเนินการของตนเอง

\subsection{ความเสี่ยงด้านข้อมูล (ระดับสูง)}
ยังไม่มีรายงานผลรายกิจกรรมเข้าระบบติดตาม (0 จาก 268~กิจกรรม) ทำให้ไม่สามารถยืนยัน
ผลสำเร็จที่ผ่านการตรวจรับได้ และการรายงานตัวชี้วัดต้องอ้างค่าที่รับมาดำเนินการแทน
หากสถานการณ์นี้คงอยู่จนสิ้นปีงบประมาณ จะกระทบต่อความน่าเชื่อถือของรายงานฉบับสมบูรณ์
และต่อการนำผลไปรวมกับตัวชี้วัดระดับมหาวิทยาลัย

\subsection{ช่องว่างเชิงการออกแบบตัวชี้วัด (ระดับปานกลาง)}
ตัวชี้วัดที่ 10 (เครือข่าย) ที่ 17 (ทรัพย์สินทางปัญญา) และที่ 38 (สถานประกอบการ)
อยู่ต่ำกว่าร้อยละ~20 เนื่องจากโครงการส่วนใหญ่ออกแบบเพื่อชุมชนและพื้นที่โครงการหลวง
มิได้ออกแบบเพื่อสถานประกอบการหรือเพื่อการจดทะเบียนทรัพย์สินทางปัญญาตั้งแต่ต้น
จึงเป็นช่องว่างเชิงโครงสร้างมากกว่าเป็นความบกพร่องของการดำเนินงาน

\subsection{ความเหลื่อมของอัตราเบิกจ่ายระหว่างหน่วยงาน (ระดับปานกลาง)}
อัตราการเบิกจ่ายกระจายตั้งแต่ร้อยละ~96 (มทร.ล้านนา น่าน) ลงมาถึงร้อยละ~0
(มทร.ล้านนา ตาก และคณะบริหารธุรกิจและศิลปศาสตร์) ความต่างในระดับนี้บ่งชี้ว่า
ปัญหาน่าจะอยู่ที่ศักยภาพด้านกระบวนการพัสดุของแต่ละหน่วยงาน มิใช่ที่ตัวเนื้องาน

\section{ข้อสรุปเชิงวิเคราะห์}

\begin{tipbox}
ปัญหาที่พบทั้งหมดเป็นเรื่อง\textbf{กระบวนการเบิกจ่าย การรายงานข้อมูล และการออกแบบตัวชี้วัด}
มิใช่เรื่องคุณภาพของงานในพื้นที่ ซึ่งยืนยันได้จากตัวชี้วัดภารกิจหลัก (ที่ 39) ที่ทำได้เกินเป้าแล้ว
ดังนั้นมาตรการที่ควรใช้จึงเป็นมาตรการด้านการบริหารจัดการและการสนับสนุนกระบวนการ
มิใช่การปรับเปลี่ยนเนื้องานของโครงการ
\end{tipbox}

\section{สรุปสถานะรายด้าน}

\begin{table}[H]
\centering
\caption{สรุปสถานะการดำเนินงานรายด้าน ณ 4 สิงหาคม 2569}
\small
\begin{tabular}{L{4.0cm}L{2.3cm}L{7.1cm}}
\toprule
\rowcolor{headerblue}
\textbf{ด้าน} & \textbf{สถานะ} & \textbf{เหตุผล}\\
\midrule
ตัวชี้วัดภารกิจหลัก (ที่ 39) & ผ่านเป้า & รับมา 62 จากเป้า 60 คิดเป็น 103\%\\
\rowcolor{rowblue}
การอนุมัติและโอนงบประมาณ & ผ่านเป้า & อนุมัติ 59 จาก 61~โครงการ โอนงบร้อยละ~98\\
การพัฒนากำลังคน (ที่ 35) & ตามแผน & รับมา 288 จากเป้า 400 คิดเป็น 72\%\\
\rowcolor{rowblue}
การเบิกจ่ายงบประมาณ & ต้องเร่ง & ร้อยละ~56 เทียบเวลาที่ผ่านไปร้อยละ~85\\
องค์ความรู้ยกระดับชุมชน (ที่ 40) & ต้องเร่ง & รับมา 45 จากเป้า 100 คิดเป็น 45\%\\
\rowcolor{rowblue}
ตัวชี้วัดเครือข่ายและเชิงพาณิชย์ & ต้องแก้เชิงโครงสร้าง & ที่ 10 18\% ที่ 17 13\% ที่ 38 8\%\\
การรายงานผลเข้าระบบ & ต้องแก้เร่งด่วน & ยังไม่มีรายงานผลรายกิจกรรม (0 จาก 268)\\
\bottomrule
\end{tabular}
\end{table}
